\begin{textsample}{NEG Dim 1 – vem\_ed\_20 – Score 9.00 – vem\_ed\_20/t100.txt}  \label{ex:f1_neg_081}
Jornada dos Projetos Colaborativos Internacionais 2023 No dia 5 de outubro, a Coordenação de Línguas e Projetos Internacionais, a Equipe de Desenvolvimento Instrucional (EDI) e a Equipe de Formação Continuada da Unidade do Ensino Superior de Graduação (Cesu) realizaram a Jornada dos Projetos Colaborativos Internacionais 2023.
O webinário gratuito trouxe diálogos sobre Intercâmbios Virtuais entre professores de Fatecs e de outras quatro instituições de ensino superior brasileiras: PUC-Campinas (São Paulo), UFF (Rio de Janeiro), Unesp (São Paulo) e Unichristus (Ceará).
A abertura contou com as participações do coordenador técnico da Cesu, Rafael Ferreira Alves; do diretor do Departamento Acadêmico-Pedagógico da Cesu, André Luiz Braun Galvão; do diretor da Faculdade de Letras e Coordenador do Departamento de Relações Externas da PUC-Campinas, Carlos Eduardo Pizzolatto; do diretor da Fatec Americana, Wladimir Costa; do diretor da Fatec Lins, Luciano Soares de Souza; da diretora da Fatec Praia Grande, Viviam Ester de Souza e de Angel Antonio Fernández Montiel, coordenador de cooperação acadêmica da Universidad Veracruzana (México).
Após uma breve explanação sobre Intercâmbios Virtuais, Internacionalização em Casa e Cidadania Global pelo coordenador dos PCIs / Cesu, Osvaldo Succi Junior, houve três painéis temáticos, cada um deles com sessões de perguntas do público. “Longevidade e manutenção de Intercâmbios Virtuais” foi o assunto debatido entre Ana Cristina Biondo Salomão (coordenadora geral do Programa Brazilian Virtual Exchange – BraVE e assessora de relações internacionais na Unesp) e Carlos Augusto Amaral Moreira (professor das Fatecs Americana, Campinas e PUC-Campinas).
Moreira desenvolve PCIs desde 2013.
Foram cerca de 30 projetos com instituições dos Estados Unidos, Holanda, França e Portugal, dentre os quais Approaching Future International Managers, realizado há 10 anos entre SUNY Ulster (EUA) e Fatec Americana.
Esse PCI está em sua vigésima primeira edição e recebeu o Prêmio Guia do Estudante / Santander Universidades de melhor parceria acadêmica em 2014.

% matched lemmas: 
\end{textsample}
